\clearpage
\fancyhead[L]{\small VON ZACH. TABULAE ABERRATIONIS ETC.}
\hypertarget{native-494-lower}{}\label{native-494-lower}
\begin{center}
\small Göttingische gelehrte Anzeigen. No. 2. pp. 18..24. 2 January 1808.\\[3pt]
\rule{.72\linewidth}{.3pt}
\end{center}

\emph{Tabulae speciales aberrationis et nutationis in ascensionem rectam et in declinationem ad supputandas stellarum fixarum positiones, sive apparentes, sive veras, una cum insigniorum 494 stellarum zodiacalium catalogo novo in specula astronomica Ernestina ad initium anni \textup{1800} constructo, cum aliis tabulis eo spectantibus, auctore} \textsc{Franc. Lib. Bar. de Zach}. Vol. I. 1806. 208 and \textsc{clvi} pp. Vol. II. 1807. 508 pp. Large octavo.

The publication of this important work, which contains the valuable fruits of many years of labour, will be most welcome to all practical astronomers, who have been looking forward to it for several years. Its main subject is the small changes in the positions of the fixed stars due to precession, aberration and nutation.%
\SourceLeaf{495}%
\space As is well known, these positions form the basis of most astronomical observations, so calculating those small changes is among the practical astronomer’s daily tasks. The celebrated author has therefore rendered practical astronomers a great service with a work that gives the most adaptable and useful formulas for those changes, works out the most reliable numerical values for the quantities on which they depend through the most careful discussion of the latest and best-tested investigations and observations, and shortens the work as much as possible with well-designed and convenient auxiliary tables. In a word, it covers everything one could wish for on the subject. Valuable as all this is, the reviewer would rate another substantial part of the work even higher, though the title does not so much as mention it. This is the catalogue of 1830 zodiacal stars, a monument more enduring than bronze to two observatories that were the pride of Germany and now, alas, both stand unused. We must devote a fuller review to a work of this importance.

The introduction takes up more than half of the first volume and begins with precession. A brief account of its original discovery and of the theoretical work of \textsc{Newton}, \textsc{d’Alembert}, \textsc{Euler}, \textsc{Lagrange} and \textsc{Laplace} is followed by a table of the variation in precession in different centuries. Then come formulas and auxiliary tables made easy to use by a device of the author’s own, for the secular change in latitude and for the part of the secular change in longitude that varies from star to star. Next come the formulas for precession in right ascension and declination, together with a table of the values of the constants from 1450 to 1950. The author found the fundamental values by making very careful comparisons of his own. He was right to use only the more recent observations for this purpose. Despite the longer interval, the observations of \textsc{Hipparchus} and \textsc{Tycho} are so much less precise that they cannot yield results anywhere near as reliable as the observations made by \textsc{Mayer} and \textsc{Bradley} with the best mural quadrants. Even \textsc{Flamsteed}’s observations fall too far short of these and would in any case have to be reduced entirely anew if anything useful were to be inferred from them. The reviewer, however, would not have counted among the reasons for excluding \textsc{Flamsteed}’s observations the claim that any slight uncertainty still remaining in the decrease of the obliquity of the ecliptic and in the proper motion of the fixed stars amounts to more for observations farther apart in time than for those closer together. By comparing 175 declinations from \textsc{Bradley} with those of \textsc{Barry}, the author derives a total annual precession of $50''1670$. Comparing the declinations of \textsc{Mayer} and \textsc{Barry} gives $49''9140$. Next, 350 right ascensions from \textsc{Mayer} compared with the author’s own determination gave $50''0718$, and finally 222 right ascensions from \textsc{Bradley} compared with the author’s gave $50''0631$. The mean of all the determinations is thus $50''0540$, the value used as the basis for all the formulas and tables in the work. We shall only touch on the investigations of the proper motions of the stars that follow. Among other things, readers will be pleased to find a completely new reduction of all \textsc{Flammsteed}’s observations of the declination of the pole star, from which a secular increase of only $3''$ follows. To calculate declination from longitude and latitude, the author uses tables with 10 decimal places here. This is certainly necessary when determining the declination from its sine, but this inconvenience is easily avoided, and one can get just as far with ordinary tables. A table giving the mean right ascension and declination of the pole star for each year from 1790 to 1820 concludes this substantial discussion.

\begingroup\rightskip=0pt plus 1em
To determine the constant of \emph{aberration}, the author re-examined all of \textsc{Bradley}’s original
\SourceLeaf{496}
observations with the greatest precision. He finds $20''232$, which differs only insignificantly from $20''255$, the result \textsc{Delambre} determined from the eclipses of Jupiter’s satellites. The author used the latter value as the basis for a general table of aberration at the end of the second volume. The table takes up only two pages, but using it requires multiplication by a sine and a secant. The author has also dealt with this minor inconvenience by a device of his own, so that using another general table, included in the introduction, requires only additions. Next come formulas and tables for aberration taking account of the Earth’s elliptical motion and for aberration due to its daily motion. Finally, there are formulas and tables for the change in aberration for a change of $1^\circ$ in the right ascension and declination of the star, from which one can estimate how long the special tables remain useful.\par\endgroup

For \emph{nutation} too, the author gives two general tables in addition to the formulas. The one at the end of the second volume takes the major and minor semiaxes of the nutation ellipse to be $10''056$ and $7''486$, following \textsc{Laplace}, and requires multiplication by a tangent. The other table, incorporated in the introduction, avoids this inconvenience. Mr \textsc{von Zach} has, however, taken the semiaxes there to be $9''648$ and $7''182$ because he thought the mass of the Moon assumed by \textsc{Laplace} had to be reduced somewhat. This section also gives instructions for calculating the part of nutation that depends on the longitude of the Sun, and formulas and tables for the change in nutation when the position of the star changes.

The rest of the introduction gives instructions for setting the transit instrument correctly, explains the catalogues and tables in this volume, and presents investigations and tables relating to the aberration of the planets, their parallax and their diameters. We cannot agree with the claim on p. 203 about the method of calculating planetary aberration that we consider preferable to all others. The claim is that this method (changing the time alone) --- gives only the part of the aberration that depends on the planet’s motion. It is understood that the position of the Earth used in the calculation must also correspond to the changed instant.

\begingroup\rightskip=0pt plus 1em
All the instructions in this substantial introduction are further illustrated by numerous well-chosen examples, so that even the least experienced reader will have no difficulty applying them. Throughout, references are given to works where one can learn more about the subjects discussed.\par\endgroup

We shall give only a brief account of the contents of the tables themselves. They give the right ascensions and declinations of the 36 principal stars according to \textsc{Maskelyne}’s latest determination, together with special aberration and nutation tables for each star. There are similar tables for the pole star for 1790, 1800, 1810 and 1820, and a list of fixed stars with nearly equal right ascensions and declinations nearly equal in magnitude but opposite in sign, for adjusting the position of the transit instrument. Next comes the most valuable part of the entire work, the right ascensions of 1830 zodiacal stars for the year 1800 determined at the Seeberg Observatory, together with the differences from \textsc{Piazzi}’s catalogue. These differences are mostly small and generally become much smaller still if one adds the correction of $3''8$ to \textsc{Piazzi}’s right ascensions. Even so, one is always entitled to prefer the determinations made with the magnificent 8-foot meridian telescope at Seeberg. Declinations are given only in minutes here, solely to identify the stars, and are missing altogether for some of them. It would be desirable for someone to take the trouble to fill this gap. The Histoire céleste would generally suffice, unless several stars in the zodiac happen to agree within one second of time in right
\SourceLeaf{497}
ascension. Many of these stars even seem to occur in \textsc{Barry}’s list of declinations immediately following, for example 921 and 963 in Virgo. To provide a sufficient number of precisely determined declinations as well, Mr \textsc{von Zach} has added a second catalogue of about 1200 zodiacal stars, containing only these declinations and the approximate right ascensions. It gives the results of observations made by \textsc{Barry} and \textsc{Henry} at the Mannheim Observatory with an 8-foot mural quadrant by \textsc{Bird}. The differences from \textsc{Piazzi}’s catalogue, also given there, show that this catalogue deserves to stand alongside his. Here too, right ascensions are missing altogether for several stars, but this is less important here, since in most cases the order alone can serve to identify the star.

The remainder of the first volume contains a general table for precession in right ascension and tables of mean refractions and their correction for the barometer and thermometer readings. The refraction tables follow the same principles as those in \textsc{Delambre}’s new solar tables and are merely reduced to the units in general use up to now. There are also tables for solar parallax and for converting sidereal time to mean and true solar time, as well as some other smaller tables.

Almost the entire \emph{second volume} is taken up by special aberration and nutation tables for 494 zodiacal stars, together with their positions for 1800 according to the author and other astronomers. We have already mentioned how valuable this aid is to practical astronomers. It remains only to note that the second volume was printed earlier than the first, and therefore differs from it on several points. All the star positions given by the author must therefore be replaced with those in the large catalogue in the volume, which have been corrected in line with \textsc{Maskelyne}’s revision of his fundamental catalogue. Furthermore, if the precessions, aberrations and nutations are to be brought into full agreement with the latest determinations in the first volume, the first must be reduced by $\frac{1}{124}$, the second increased by $\frac{1}{78}$, and the third increased by $\frac{1}{92}$.

The second volume concludes with the general tables for aberration and nutation already mentioned, and a table for the longitude of the ascending node of the Moon’s orbit based on \textsc{Bürg}’s lunar tables.
